% Standalone problem document, generated from the adjacent README.md.
% Compile with XeLaTeX. Mathematical content is maintained in Markdown.
\documentclass[11pt,a4paper]{article}
\usepackage[fontsize=10.5pt]{scrextend}
\usepackage[margin=22mm,headheight=15pt,headsep=9mm,footskip=11mm]{geometry}
\usepackage{fontspec}
\setmainfont{texgyrepagella}[Extension=.otf,UprightFont=*-regular,BoldFont=*-bold,ItalicFont=*-italic,BoldItalicFont=*-bolditalic]
\setsansfont{texgyreheros}[Extension=.otf,UprightFont=*-regular,BoldFont=*-bold,ItalicFont=*-italic,BoldItalicFont=*-bolditalic]
\setmonofont{texgyrecursor}[Extension=.otf,UprightFont=*-regular,BoldFont=*-bold,ItalicFont=*-italic,BoldItalicFont=*-bolditalic,Scale=MatchLowercase]
\usepackage{amsmath,amssymb}
\usepackage{unicode-math}
\setmathfont{texgyrepagella-math.otf}
\usepackage{xcolor,microtype,enumitem,needspace,fancyhdr}
\usepackage{bookmark}
\definecolor{ink}{HTML}{17344B}
\definecolor{accent}{HTML}{197D80}
\definecolor{muted}{HTML}{596773}
\hypersetup{colorlinks=true,linkcolor=accent,urlcolor=accent,pdftitle={TR-30 - Dimension
of tensor loci with prescribed minimum border
subrank},pdfauthor={Open Problems in Numerical Linear Algebra}}
\urlstyle{same}
\setlength{\parindent}{0pt}
\setlength{\parskip}{5pt plus 1pt minus 1pt}
\linespread{1.04}
\setlength{\emergencystretch}{3em}
\widowpenalty=10000
\clubpenalty=10000
\displaywidowpenalty=10000
\allowdisplaybreaks[1]
\setlist{nosep,leftmargin=1.5em}
\providecommand{\tightlist}{\setlength{\itemsep}{0pt}\setlength{\parskip}{0pt}}
\providecommand{\pandocbounded}[1]{#1}
\pagestyle{fancy}
\fancyhf{}
\fancyhead[L]{\sffamily\footnotesize\color{muted}OPEN PROBLEMS IN NUMERICAL LINEAR ALGEBRA}
\fancyhead[R]{\sffamily\bfseries\footnotesize\color{accent}TR-30}
\fancyfoot[L]{\parbox[b]{0.9\textwidth}{\sffamily\fontsize{8}{10}\selectfont\color{muted}Editorial ratings; a bounded search cannot certify that a problem remains unsolved.\\Literature check: 2026-09-10}}
\fancyfoot[R]{\sffamily\footnotesize\color{muted}\thepage}
\renewcommand{\headrulewidth}{0.3pt}
\renewcommand{\headrule}{\hbox to\headwidth{\color{accent}\leaders\hrule height \headrulewidth\hfill}}
\makeatletter
\renewcommand{\section}{\Needspace{5\baselineskip}\@startsection{section}{1}{0pt}{12pt plus 2pt minus 2pt}{4pt}{\sffamily\bfseries\large\color{ink}}}
\renewcommand{\subsection}{\Needspace{4\baselineskip}\@startsection{subsection}{2}{0pt}{9pt plus 1pt minus 1pt}{3pt}{\sffamily\bfseries\normalsize\color{ink}}}
\makeatother
\setcounter{secnumdepth}{0}
\begin{document}
{\sffamily\small\color{muted}Tensor computations\par}
\vspace{3pt}
{\raggedright\hyphenpenalty=10000\exhyphenpenalty=10000\sffamily\bfseries\fontsize{21}{24}\selectfont\color{ink}Dimension
of tensor loci with prescribed minimum border subrank\par}
\vspace{7pt}
{\sffamily\small\textbf{Difficulty:} challenging\quad\textbf{Importance:} interesting
to specialist\par}
{\sffamily\small\bfseries\color{accent}Status: Partially resolved\par}
\vspace{4pt}
{\color{accent}\hrule height 0.8pt}
\vspace{6pt}
\textbf{Rating rationale:} Exact dimensions beyond the generic range
require substantial control of tensor degenerations and stronger
arguments than current bounds. The immediate audience is specialists in
the geometry of tensor subrank.

\subsection{Problem statement}\label{problem-statement}

Let \(\mathbb F\) be an algebraically closed field, let \(k\geq3\), and
fix positive integers \(n_1,\ldots,n_k\) and \(1\leq r\leq\min_jn_j\).
For \(T\in V=\mathbb F^{n_1}\otimes\cdots\otimes\mathbb F^{n_k}\), say
that \(\underline Q(T)\geq r\) when \[
\sum_{a=1}^r e_a^{\otimes k}\in
\overline{\{(L_1\otimes\cdots\otimes L_k)T:
L_j\in\operatorname{Hom}(\mathbb F^{n_j},\mathbb F^r)\}}^{\,\mathrm{Zar}}.
\] Thus \(\underline Q(T)\) is border subrank, defined through
degeneration to a diagonal unit tensor.

Determine, for every such field and tuple of integers, the dimension \[
\dim X_r^{(n_1,\ldots,n_k)},\qquad
X_r^{(n_1,\ldots,n_k)}
=\{T\in V:\underline Q(T)\geq r\}.
\] The dimension of this constructible locus means the Krull dimension
of its Zariski closure in the affine tensor space. This convention does
not assert that the locus itself is closed. The requested output is its
exact dimension as a function of the parameters; determining only an
asymptotic order or only the generic subrank does not complete the
question.

\subsection{Why it matters}\label{why-it-matters}

Border subrank measures how many independent diagonal computations a
tensor can approximate after transformations of its modes. The dimension
asks how large the families with a specified computational capability
are. It is a geometric complement to the better-understood
low-border-rank varieties used in decomposition problems.

\newpage
\subsection{References and status
check}\label{references-and-status-check}

\begin{enumerate}
\def\labelenumi{\arabic{enumi}.}
\item
  A. Oneto and E. Ventura, \emph{Ranks of tensors: geometry and
  applications}, \href{https://doi.org/10.1007/s40574-025-00472-9}{2025
  survey}, Question 14 in §4.4.
\item
  C.-Y. Chang, \emph{Maximal border subrank tensors},
  \href{https://arxiv.org/abs/2208.04281}{arXiv:2208.04281};
  \href{https://doi.org/10.1080/03081087.2024.2352456}{Linear and
  Multilinear Algebra \textbf{73} (2025), 525--535}, the dimension
  lower-bound theorem.
\item
  B. Biaggi, C.-Y. Chang, J. Draisma, and F. Rupniewski, \emph{Border
  subrank via a generalised Hilbert--Mumford criterion},
  \href{https://arxiv.org/html/2402.10674v2}{arXiv:2402.10674v2}, §1
  (arbitrary algebraically closed fields), Proposition 1
  (constructibility), and Theorem 3 (dimension upper bound); Advances in
  Mathematics \textbf{461} (2025), 110077.
\item
  P. Pielasa, M. Šafránek, and A. Shatsila, \emph{Exact values of
  generic subrank},
  \href{https://arxiv.org/pdf/2408.07550}{arXiv:2408.07550}, Theorem 3.7
  (all infinite fields).
\item
  J. Draisma, \emph{(Border) subranks of tensors},
  \href{https://www.combinatorial-synergies.de/activities/2026-09_AnnualConference/}{September
  2026 conference abstract}, ``Ordinary border subrank'' bullet; an
  announcement of a restricted result.
\end{enumerate}

\subsubsection{Status check ---
2026-09-10}\label{status-check-2026-09-10}

Checked the 2025 source question, Biaggi--Chang--Draisma--Rupniewski v2,
and targeted dimension searches through 2026. There is a proved exact
range: by Pielasa--Šafránek--Shatsila, Theorem 3.7, let
\(g=\min\{n_1,\ldots,n_k,\lfloor(\sum_i n_i-k+1)^{1/(k-1)}\rfloor\}\).
Generic tensors have ordinary subrank \(g\) over every infinite field.
Since border subrank is at least ordinary subrank, \(X_r\) contains a
dense open subset and has dimension \(\prod_i n_i\) whenever \(r\le g\).
This deduction resolves that range, not the remaining dimensions.
Draisma's abstract for the September 14--16, 2026 conference announces
sharp bounds for equal-format tensors of maximal border subrank; the
abstract supplies no proof or all-format formula. No complete resolution
was located.
\end{document}
